\section{A finite semidefinite lift}\label{sec:lift}

We now replace the angular family of pencils by one finite matrix condition.
Evaluation on an angular slice will provide a feasible lift for each point of
$K$.  Conversely, the functions supplied by Lemma~\ref{lem:matching} will
turn positivity of the finite matrix into the inequalities that characterize
the slice, including its singular boundary.

\subsection{The matrix condition}

Let $\mathcal T_j$ be the real vector space of trigonometric polynomials of
degree at most $j$, with basis
\[
  1,\cos\theta,\sin\theta,\ldots,\cos(j\theta),\sin(j\theta).
\]
Let $\mathcal E$ be the space of real functions of
$(\theta,\mathcal X,\mathcal Z,\eta,\xi)$ that are linear homogeneous in
the $22$ coordinates of
$(\mathcal X,\mathcal Z,\eta,\xi)\in
\mathbb S^4\times\mathbb S^4\times\mathbb R^2$, with coefficients in
$\mathcal T_7$.  Thus $\dim\mathcal E=22\cdot15=330$.
Our lifting variable is a real linear functional
$\Lambda:\mathcal E\to\mathbb R$; on vectors and matrices we apply it
entrywise.  For an output point $(X,Z,y)$, impose
\begin{equation}\label{eq:output}
  \Lambda(\mathcal X)=X,\qquad
  \Lambda(\mathcal Z)=Z,\qquad
  \Lambda\bigl(Y_\theta(\eta,\xi)\bigr)=y.
\end{equation}
We also require
\begin{equation}\label{eq:moment}
  \Lambda\bigl(q(\theta)^T
  L_\theta(\mathcal X,\mathcal Z,\eta,\xi)q(\theta)\bigr)\geq0
  \qquad\text{for every }q\in(\mathcal T_2)^{20}.
\end{equation}
The expression to which $\Lambda$ is applied lies in $\mathcal E$:
the entries of $L_\theta$ have trigonometric degree at most three, so its
product with two entries of $q$ has degree at most seven.

Condition~\eqref{eq:moment} is a single finite semidefinite condition,
using the finite test-space construction familiar from moment
relaxations~\cite{lasserre2001,lasserre2009}.
Indeed, set $(f_1,\ldots,f_5)=
(1,\cos\theta,\sin\theta,\cos2\theta,\sin2\theta)$ and define
$H(\Lambda)\in\mathbb S^{100}$ by
\begin{equation}\label{eq:moment-matrix}
  H(\Lambda)_{(i,a),(j,b)}
    =\Lambda\bigl(f_a f_b (L_\theta)_{ij}\bigr),
  \qquad 1\leq i,j\leq20,\quad 1\leq a,b\leq5.
\end{equation}
If $q_i=\sum_a t_{i,a}f_a$, the left side of~\eqref{eq:moment} is
$t^TH(\Lambda)t$.  Hence~\eqref{eq:moment} is equivalent to
$H(\Lambda)\succeq0$.  Every entry of this real symmetric matrix is linear
homogeneous in the $330$ coordinates of $\Lambda$.

\begin{samepage}
\begin{proposition}\label{prop:lift}
For $(X,Z,y)\in\mathbb S^4\times\mathbb S^4\times\mathbb R^3$, the
following are equivalent:
\begin{enumerate}
  \item $(X,Z,y)\in K$;
  \item there is a real linear functional $\Lambda:\mathcal E\to\mathbb R$
  satisfying~\eqref{eq:output} and $H(\Lambda)\succeq0$.
\end{enumerate}
\end{proposition}
\end{samepage}

\begin{proof}
Suppose first that $(X,Z,y)\in K$.  Choose
$\theta_0,\eta_0,\xi_0$ with
$y=Y_{\theta_0}(\eta_0,\xi_0)$, and let $\Lambda$ be evaluation at
$(\theta_0,X,Z,\eta_0,\xi_0)$.  The output equations hold, and
Lemma~\ref{lem:slice} gives
$L_{\theta_0}(X,Z,\eta_0,\xi_0)\succeq0$, which implies
\eqref{eq:moment}.  This evaluation functional is a finite feasible
witness even when $X$ is singular.

For the converse, suppose that $\Lambda$ satisfies the stated conditions.
Taking constant $q$ supported on one of the four top blocks of $L_\theta$
in~\eqref{eq:moment} gives $X\succeq0$.
We next obtain a family of scalar inequalities from which the remaining
slice conditions will follow.

Fix $u\in\mathbb R^4$, $D\in\mathbb S^4$, and
$z\in\mathbb R^3$ with $(z_1,z_2)\neq(0,0)$, and apply
Lemma~\ref{lem:matching} with
$S=J_uDJ_u^T\in\mathbb S^3$.  Neither $D$ nor $S$ is required to be
positive semidefinite.  Let $b_1(\theta),\ldots,b_4(\theta)$ be
the columns of the resulting real matrix $\widetilde C(\theta)$.
Choose the five four-dimensional blocks of $q$ to be
\[
  q(\theta)=
  \begin{pmatrix}
    -DJ_u^Tb_1(\theta)\\
    \vdots\\
    -DJ_u^Tb_4(\theta)\\
    u
  \end{pmatrix}.
\]
The degree bound in Lemma~\ref{lem:matching} puts this vector in
$(\mathcal T_2)^{20}$.  Write $r_j$ for the columns of
$\widetilde F_\theta(\eta,\xi)$, as in the definition of $L_\theta$.
Using $B(r_j)u=J_u^Tr_j$, the top blocks and the cross terms of
$q^TL_\theta q$ are respectively
\begin{align*}
  \sum_{j=1}^4 b_j^T J_uD\mathcal X DJ_u^Tb_j
    &=\operatorname{tr}\bigl(J_uD\mathcal X DJ_u^TQ(z)\bigr),\\
  -2\sum_{j=1}^4 b_j^TJ_uDJ_u^Tr_j
    &=-2\langle F_\theta(\eta,\xi),C(\theta)\rangle_S\\
    &=-2\operatorname{tr}\bigl(J_uDJ_u^T
        Q(Y_\theta(\eta,\xi),z)\bigr).
\end{align*}
These identities hold before applying $\Lambda$.  The first identity
removes the angular dependence from the coefficient of $\mathcal X$;
the last expression is linear in $Y_\theta(\eta,\xi)$.
Consequently, applying~\eqref{eq:output} and~\eqref{eq:moment} gives
\begin{equation}\label{eq:certificate}
  u^TZu+
  \operatorname{tr}\bigl(J_uDXDJ_u^TQ(z)\bigr)
  -2\operatorname{tr}\bigl(J_uDJ_u^TQ(y,z)\bigr)\geq0.
\end{equation}
For fixed $u,D$, the left side is a polynomial in $z$, so the inequality
extends by continuity to every $z\in\mathbb R^3$.  Only this scalar
inequality is extended; no continuity of the functions $C$ is needed.
Thus~\eqref{eq:certificate} holds for all $u\in\mathbb R^4$,
$D\in\mathbb S^4$, and $z\in\mathbb R^3$.

If $X$ were positive definite, setting $z=y$ and $D=X^{-1}$ in
\eqref{eq:certificate} would give $Z\succeq\Phi_y(X^{-1})$.
For singular $X$, the same family of tests also supplies the necessary
range condition.  Fix slice parameters
$y=Y_{\theta_0}(\eta_0,\xi_0)$ and, for the rest of the proof, let
$r_j$ denote the columns of
$\widetilde F_{\theta_0}(\eta_0,\xi_0)$ and set $B_j=B(r_j)$.
For $v\in\ker X$, take $D=t vv^T$ and $z=y$ in
\eqref{eq:certificate}, where $t\in\mathbb R$ is arbitrary.  Since
$DXD=0$, this gives
\[
  u^TZu-2t\,(J_uv)^TQ(y)(J_uv)\geq0
  \qquad(t\in\mathbb R).
\]
Its coefficient of $t$ must vanish.  The Gram factorization of $Q(y)$
therefore yields
\[
  0=(J_uv)^TQ(y)(J_uv)
    =\sum_{j=1}^4(v^TB_ju)^2.
\]
As $u$ and $v\in\ker X$ are arbitrary, it follows that
\begin{equation}\label{eq:range}
  \operatorname{range}B_j\subseteq(\ker X)^\perp
    =\operatorname{range}X\qquad(1\leq j\leq4).
\end{equation}

Let $X^\dagger$ be the symmetric matrix that inverts the positive
eigenvalues of $X$ and is zero on $\ker X$.  In particular,
$X^\dagger XX^\dagger=X^\dagger$.
Taking $D=X^\dagger$ and $z=y$ in~\eqref{eq:certificate}, and using
\eqref{eq:Phi-quadratic} and~\eqref{eq:Phi-factor}, gives
\begin{equation}\label{eq:singular-schur}
  Z\succeq\Phi_y(X^\dagger)
    =\sum_{j=1}^4B_j^TX^\dagger B_j.
\end{equation}
We now verify positivity of the slice pencil directly by completing
a square using~\eqref{eq:range} and~\eqref{eq:singular-schur}.  For
$a_1,\ldots,a_4,u\in\mathbb R^4$, its quadratic form equals
\begin{align*}
  &\sum_{j=1}^4a_j^TXa_j
    +2\sum_{j=1}^4a_j^TB_ju+u^TZu\\
  &\quad=\sum_{j=1}^4
    (a_j+X^\dagger B_ju)^TX(a_j+X^\dagger B_ju)
    +u^T\left(Z-\sum_{j=1}^4B_j^TX^\dagger B_j\right)u
    \geq0.
\end{align*}
Here the cross terms agree because $XX^\dagger B_j=B_j$ by
\eqref{eq:range}.  Thus
$L_{\theta_0}(X,Z,\eta_0,\xi_0)\succeq0$, and
Lemma~\ref{lem:slice} gives $(X,Z,y)\in K$.
\end{proof}

The converse has used an arbitrary feasible functional $\Lambda$.
In particular, the proof requires no representation of $\Lambda$ by a
measure and no closure operation on the projected feasible set.

\subsection{The size of the lift}

\begin{proof}[Proof of Theorem~\ref{thm:main}]
Choose the displayed trigonometric basis of $\mathcal T_7$ and the
upper-triangular coordinate functions of $\mathcal X$ and $\mathcal Z$,
together with $\eta,\xi$, to specify the $330$ coordinates of $\Lambda$.
The output equations~\eqref{eq:output} prescribe precisely the following
$23$ distinct coordinates:
\begin{gather*}
  \Lambda(\mathcal X_{ij}),\quad
  \Lambda(\mathcal Z_{ij})\qquad(1\leq i\leq j\leq4),\\
  \Lambda(\eta\cos\theta),\quad
  \Lambda(\eta\sin\theta),\quad
  \Lambda(\xi).
\end{gather*}
They are independent members of the chosen coordinate system on
$\mathcal E^*$.  Substitute the corresponding entries of $X,Z,y$ for
these coordinates in~\eqref{eq:moment-matrix}, and use the remaining
$330-23=307$ coordinates as auxiliary variables.  The result is a fixed
real symmetric $100\times100$ matrix pencil, homogeneous in the output
and auxiliary variables.  Proposition~\ref{prop:lift} proves that its
projection is exactly $K$.  Hence the requested representation has
$N=100$, $m=307$, and zero constant term.
\end{proof}
